\section*{Bab \ref{ch:b1:populations} --- Populasi dan Demografi}

\begin{solution}{exo:b1:populations:1}
\omterm{def:b1:populations:population}{Populasi}: individu satu spesies di satu tempat dan waktu, yang sanggup
berbiak silang. Kerapatan: individu tiap satuan luas atau volume.
Sebaran: pola ruangnya --- mengelompok (sekawanan, serumpun jelatang),
seragam (burung gannet bersarang sejauh satu patukan, semak kreosot),
acak (jombang di sebuah halaman rumput).
\end{solution}

\begin{solution}{exo:b1:populations:2}
$l_x$: fraksi sebuah kohort yang hidup pada umur $x$. $m_x$: anak
betina tiap betina berumur $x$. $R_0 = \sum l_x m_x$: anak betina tiap
betina yang baru lahir sepanjang hidupnya. $T = \sum x l_x m_x / R_0$:
umur purata induknya. $R_0 = 1$: tiap betina menggantikan dirinya;
\omterm{def:b1:populations:population}{populasinya} mantap.
\end{solution}

\begin{solution}{exo:b1:populations:3}
Jenis III: sekitar 15 dari 1000, \qty{1.5}{\%}. Jenis I: sekitar 700 dari
1000, \qty{70}{\%}.
\end{solution}

\begin{solution}{exo:b1:populations:4}
$\dd N/\dd t = rN$ dan $\dd N/\dd t = rN(1 - N/K)$. $r$ adalah laju
pertumbuhan tiap kepalanya ketika sumber dayanya tak terbatas
(kelahiran dikurangi kematian tiap individu tiap satuan waktu); $K$
adalah \omterm{thm:b1:populations:logistic}{daya dukungnya}, yaitu ukuran ketika pertumbuhannya berhenti.
\end{solution}

\begin{solution}{exo:b1:populations:5}
$r = \ln(1000)/5 = \qty{1.38}{h^{-1}}$; lipat duanya $\ln 2/1.38 =
\qty{0.5}{h}$.
\end{solution}

\begin{solution}{exo:b1:populations:6}
$l_x m_x = 0, 0.5, 0.6, 0.2, 0$: $R_0 = 1.3$. $T = (1\times 0.5 +
2\times 0.6 + 3\times 0.2)/1.3 = 2.3/1.3 = 1.77$. $r = \ln 1.3/1.77 =
\qty{0.15}{}$ tiap satuan waktu: ia tumbuh.
\end{solution}

\begin{solution}{exo:b1:populations:7}
$N = 50\times 70/7 = 500$. Hanya 40 ikan bertanda yang tersisa: $N = 40\times
70/7 = 400$; taksiran pertamanya terlalu tinggi seperempat (fraksi
bertanda \omterm{def:b1:populations:population}{populasinya} lebih kecil daripada yang dianggap).
\end{solution}

\begin{solution}{exo:b1:populations:8}
$0.2\times 100\times 0.9 = 18$; $0.2\times 500\times 0.5 = 50$;
$0.2\times 900\times 0.1 = 18$ tiap tahun. Maksimum $rK/4 = 50$ setahun, pada
$N = 500$.
\end{solution}

\begin{solution}{exo:b1:populations:9}
Jombang: $r$ (ribuan biji, setahunan). Gajah: $K$ (satu anak dalam
empat tahun, umur panjang, pengasuhan). Ikan kod: $r$ (jutaan telur,
tanpa asuhan). Gorila: $K$ (satu bayi dalam empat tahun, berpuluh
tahun hidup). Lalat rumah: $r$ (ratusan telur, matang dalam sepuluh
hari).
\end{solution}

\begin{solution}{exo:b1:populations:10}
$r = \ln(29/10)/2 = \qty{0.53}{h^{-1}}$; $K \approx 660$. Pada \qty{8}{h}:
$660/(1 + 65e^{-4.24}) = 660/(1 + 0.94) = 340$, lawan 351 yang dihitung.
Tercepat pada $K/2 = 330$, yang dicapai pada $t = \ln 65/0.53 = \qty{7.9}{h}$,
dengan $\dd N/\dd t = rK/4 = 87\times 10^5$ \omterm{def:b1:cell-unit-of-life:cell}{sel} tiap mililiter tiap
jam.
\end{solution}

\begin{solution}{exo:b1:populations:11}
Pengaturannya menuntut laju tiap kepalanya turun seiring $N$ naik,
sehingga $N$ kembali menuju sebuah kesetimbangan sesudah sebuah
gangguan; sebuah faktor yang membunuh fraksi yang tetap mengubah $N$
tetapi bukan kebergantungan lajunya pada $N$, sehingga sesudahnya
\omterm{def:b1:populations:population}{populasinya} melanjutkan pertumbuhan yang sama. Cuaca membatasi
serangganya dengan menyerang cukup sering, dan secara acak, sehingga
\omterm{def:b1:populations:population}{populasinya} terpukul mundur sebelum ia sanggup mendekati langit-langit
mana pun: sebuah \omterm{def:b1:populations:population}{populasi} tak teratur yang ditahan rendah oleh
gangguan berulang, yang ukuran puratanya bergantung pada kekerapan
musim yang buruk alih-alih pada sebuah \omterm{thm:b1:populations:logistic}{daya dukung}.
\end{solution}

\begin{solution}{exo:b1:populations:12}
Dua \omterm{def:b1:populations:population}{populasi} berukuran sama, satu dengan alas muda yang lebar dan satu
dengan kelas umur yang merata, memiliki masa depan yang berbeda: yang
pertama akan tumbuh berpuluh tahun bahkan bila $m_x$-nya turun ke
penggantian sekarang, sebab banyak yang mudanya belum berbiak
(momentum demografi); yang kedua tidak. $N$ merekam masa lalunya;
sebaran umurnya dan jadwal $m_x$ di atasnya itulah yang dikerjai
persamaannya, dan piramidanya meramalkan ukuran generasi berikutnya
sedangkan $N$ tak mengatakan apa pun tentangnya.
\end{solution}

\begin{solution}{pb:b1:populations:1}
\textbf{1.} $r = \ln(29/10)/2 = \qty{0.53}{h^{-1}}$.
\textbf{2.} $\ln 2/0.53 = \qty{1.3}{h}$.
\textbf{3.} $K \approx 660\times 10^5 = \qty{6.6e7}{}$ \omterm{def:b1:cell-unit-of-life:cell}{sel} tiap
mililiter.
\textbf{4.} $N(t) = 660/(1 + 65e^{-0.53t})$: pada 4 h, $660/(1 + 7.8) =
75$ (71 yang dihitung); pada 8 h, 340 (351); pada 12 h, $660/(1 + 0.112) = 594$
(595).
\textbf{5.} Pada $K/2 = 330$, yang dicapai pada $t = \ln 65/0.53 = \qty{7.9}{h}$;
laju maksimumnya $rK/4 = 0.53\times 660/4 = 87\times 10^5$ \omterm{def:b1:cell-unit-of-life:cell}{sel} tiap
mililiter tiap jam.
\textbf{6.} $10e^{0.53\times 12} = 10\times 578 = 5780$: sepuluh kali
595 yang teramati.
\textbf{7.} Habisnya gulanya (atau oksigennya); bertimbunnya
etanol dan asam yang meracuni \omterm{def:b1:cell-unit-of-life:cell}{selnya}.
\textbf{8.} $q_0 = 0.40$; $q_4 = 1 - 0.34/0.40 = 0.15$: kematian awal
yang berat lalu sebuah laju yang lebih mantap --- antara jenis II dan
III, sebagaimana pada sebagian besar mamalia besar di alam liar.
\textbf{9.} $l_x m_x = 0, 0, 0.312, 0.368, 0.320, 0.272, 0.182, 0.080,
0.018, 0$: $R_0 = 1.55$.
\textbf{10.} $\sum x l_x m_x = 0.624 + 1.104 + 1.280 + 1.360 + 1.092 +
0.560 + 0.144 = 6.16$; $T = 6.16/1.55 = \qty{4.0}{yr}$.
\textbf{11.} $r = \ln 1.55/4.0 = \qty{0.11}{yr^{-1}}$; lipat dua dalam
\qty{6.3}{yr}.
\textbf{12.} $200e^{1.1} = 600$ rusa.
\textbf{13.} $N(10) = 600/(1 + 2e^{-1.1}) = 600/1.67 = 360$; melampaui 500
ketika $(K - N_0)/N_0\,e^{-rt} = 0.2$, yakni $e^{-rt} = 0.1$, $t = \ln
10/0.11 = 21$ tahun.
\textbf{14.} $120\times 150/18 = 1000$ ikan perch.
\textbf{15.} 100 bertanda: $100\times 150/18 = 833$; taksiran pertamanya
terlalu tinggi (lebih sedikit ikan bertanda yang berkeliaran daripada
yang dianggap).
\textbf{16.} $rK/4 = 100$ setahun, pada $N = 500$.
\textbf{17.} $T_2 = \ln 2/0.4 = \qty{1.7}{yr}$.
\textbf{18.} Tangkapan $hN$ dengan $h = 0.2$: $rN(1 - N/K) = hN$ memberi $N^* =
K(1 - h/r) = 1000\times 0.5 = 500$ dan sebuah tangkapan \num{100} setahun
--- di bawah hasil lestari maksimumnya, tetapi sebuah fraksi menyetel
dirinya pada stoknya: bila stoknya terparuh, tangkapannya pun, dan
\omterm{def:b1:populations:population}{populasinya} kembali.
\textbf{19.} Pada 500: $0.4\times 500\times 0.5 = 100 < 120$: \omterm{def:b1:populations:population}{populasinya}
menyusut. Pada 300: $0.4\times 300\times 0.7 = 84 < 120$: menyusut lebih
cepat --- begitu berada di bawah $K/2$ sebuah panen tetap di atas hasil
maksimumnya menyeretnya ke kepunahan.
\textbf{20.} Pada $K/2$ kelebihannya berada pada maksimumnya tetapi
kekeliruan apa pun --- sebuah \omterm{def:b1:populations:population}{populasi} yang tertaksir terlalu tinggi
seperlima, sebuah tahun buruk yang memotong $r$ --- mengubah sebuah
pengambilan yang lestari menjadi sebuah penyusutan, dan di bawah $K/2$
kelebihannya mengecil seiring \omterm{def:b1:populations:population}{populasinya} mengecil: kesetimbangannya
tak mantap terhadap panen berlebih. Memanen agak di atas $K/2$
meninggalkan sebuah marjin.
\textbf{21.} $R_0$ terparuh menjadi $0.78 < 1$: kawanannya menyusut
sampai ia jatuh di bawah 400 dan $m_x$-nya pulih --- sebuah faktor yang
bergantung kerapatan, yang mengatur kawanannya dekat 400.
\textbf{22.} Tidak: ia menyingkirkan \qty{40}{\%} berapa pun ukurannya
lalu meninggalkan laju tiap kepalanya tak berubah; pada tahun
sesudahnya kawanannya tumbuh kembali pada $r$ yang sama dari sebuah
awal yang lebih rendah --- sebuah gangguan, bukan sebuah pengaturan.
\textbf{23.} $0.11N(1 - N/600) = 30$: $N^2 - 600N + 163\,600 = 0$,
$N = 300\pm\sqrt{90\,000 - 163\,600}$ --- tak ada penyelesaian nyata: 30
setahun melampaui kelebihan maksimumnya $rK/4 = 16.5$, dan serigalanya
menyeret kawanannya ke kepunahan. Dengan 15 setahun: $N = 300\pm\sqrt{90\,000 -
81\,800} = 300\pm 91$, yakni 209 dan 391; yang atas mantap (di atasnya
kelebihannya kurang dari 15 dan kawanannya jatuh kembali, di bawahnya
lebih dan ia naik), yang bawah tak mantap.
\textbf{24.} Jumlah lynxnya menanggapi makanannya dengan sebuah
tundaan: naiknya terwelu menaikkan kelahiran dan \omterm{def:b1:populations:lifetable}{kesintasan} lynxnya
sepanjang satu atau dua tahun berikutnya, dan runtuhnya terwelu
membuat lynxnya kelaparan hanya sesudah \omterm{def:b1:lipids:triglyceride}{lemaknya} dan anak semusimnya
habis.
\textbf{25.} Khamir: $K = \qty{6.6e7}{}$ \omterm{def:b1:cell-unit-of-life:cell}{sel} tiap mililiter, $r =
\qty{0.53}{h^{-1}}$. Rusa: $R_0 = 1.55$, $T = \qty{4.0}{yr}$, $r =
\qty{0.11}{yr^{-1}}$.
\end{solution}
